% =====================================================================
% Track A (Closed-World, from scratch): context, Phases~1--4, SSL runs.
%
% \input{}-ed from report/main.tex.  Phase~3 V-JEPA + Phase~4 + appendices
% were merged from the former monolithic main.tex; legacy DINO Phase~3 uses
% \texttt{sec:phase3dino:*} labels to avoid clashing with V-JEPA \texttt{sec:phase3:*}.
% =====================================================================

% Set Track A accent color (red stripe + red headers)
\renewcommand{\trackcolor}{trackred}
\renewcommand{\trackname}{Track~A}

\section{Context}
\label{sec:context}

\subsection{Competition rules}

The Kaggle competition \emph{Can a model predict the future?} provides
short video clips from the Something-Something~V2 dataset and asks
participants to predict the action that will occur next, given only the
beginning of each clip. Track~A (\emph{Closed World}) is the focus of this
work: \textbf{models must be trained from scratch using only the provided
data}; no external dataset and no pretrained weights are allowed.
Submissions are evaluated by Top-1 classification accuracy (with Top-5 as
a secondary metric) on a held-out test set whose public/private split is
hidden from participants.

\subsection{Starting point and observed problems}

Before the work in this report, the supervised Track-A pipeline was a
TSM-ResNet50 trained with Adam, label smoothing, RandAugment-T temporal
augmentations, FrameMixUp, an 80/20 internal split of the training data
for validation, and a cosine learning-rate schedule. Its best public
Kaggle scores were:

\begin{center}
\begin{tabular}{lcc}
\toprule
\textbf{Submission} & \textbf{Internal val Top-1} & \textbf{Public Top-1} \\
\midrule
Initial run (epoch 29) & 0.4582 & 0.3207 \\
Resumed run (epoch 60) & 0.5666 & 0.3584 \\
\bottomrule
\end{tabular}
\end{center}

The roughly 21~percentage-point gap between internal validation and the
public leaderboard pointed at a collection of issues, the most important
of which were identified in a project review:

\begin{enumerate}[leftmargin=*]
\item \textbf{Leaky validation split}. The 80/20 split was carved out of
      \file{data/train/} rather than using the official held-out
      \file{data/val/} folder. Validation samples were therefore drawn
      from the same recording sessions and demographics as training,
      producing optimistically biased numbers.
\item \textbf{Suboptimal initialisation}. The ResNet-50 backbone was
      initialised with Xavier weights, which is well-suited to networks
      with sigmoid/tanh activations but is documented to under-perform
      Kaiming He initialisation on deep ReLU networks trained from
      scratch.
\item \textbf{No mixed-precision training}. Forward and backward passes
      ran in full FP32, leaving \(\approx 1.5{-}2\times\) of GPU
      throughput on the table on Ampere or newer hardware.
\item \textbf{Optimizer / scheduler / scaler state lost on resume}. Only
      the model weights and the epoch counter were checkpointed, so
      restarting a long run forced the optimizer (Adam moments) and the
      cosine LR schedule to restart from scratch as well.
\item \textbf{Missing class~27}. The ``50 classes'' described in the
      competition page is in practice 33 unique class indices on disk
      with class~27 entirely absent. The submission pipeline did nothing
      to prevent the network from emitting this never-trained label,
      meaning \(\approx 3\%\) of test predictions could land on it.
\item \textbf{No EMA, no TTA, no SSL pretraining.} The pipeline used
      none of the well-known accuracy-boosting techniques that are
      standard for video classification on small datasets.
\end{enumerate}

The remainder of this document describes the changes that addressed each
of these issues, in the order they were introduced.

% =====================================================================
\section{Training Pipeline Foundations}
\label{sec:phase1}

Phase~1 comprises five low-risk, high-return changes. All of them were
introduced behind explicit configuration switches so that historical runs
can still be reproduced bit-for-bit by setting the new switches to their
``off'' defaults.

\subsection{Official validation split}
\label{sec:phase1:val}

\paragraph{Logic.} A train/val split that overlaps with the test
distribution understates the generalisation gap. Switching to the
official \file{data/val/} folder makes the validation score a faithful
proxy for the public leaderboard score and therefore drives early
stopping, model selection and hyperparameter choices in the right
direction.

\paragraph{Implementation.} A new boolean
\cfgkey{dataset.use\_official\_val} (default \code{false}) was added to
\file{configs/data/default.yaml}. When set to \code{true}, the trainer
bypasses the internal 80/20 split and instead enumerates every video
under \file{data/val/}.

\begin{lstlisting}[style=py]
use_official_val = bool(cfg.dataset.get("use_official_val", False))
if use_official_val:
    val_dir_for_val = Path(cfg.dataset.val_dir).resolve()
    val_samples = collect_video_samples(val_dir_for_val)
    train_samples = all_samples
    print(f"[data] use_official_val=true: train={len(train_samples)}, "
          f"val={len(val_samples)} (from {val_dir_for_val})")
else:
    train_samples, val_samples = split_train_val(
        all_samples, val_ratio=float(cfg.dataset.val_ratio),
        seed=int(cfg.dataset.seed))
\end{lstlisting}

\subsection{Kaiming He initialisation}
\label{sec:phase1:init}

\paragraph{Logic.} For a Conv$\to$BN$\to$ReLU stack trained from scratch,
the variance of pre-activations grows or shrinks geometrically with depth
unless the initial Conv weights are scaled with the Kaiming He recipe
\cite[Sec.~2.2]{HeICCV2015}:
\[
\mathrm{Var}(W_{i,j}) = \frac{2}{\text{fan\_out}}.
\]
Xavier initialisation, which uses \(2/(\text{fan\_in}+\text{fan\_out})\),
under-shoots this variance by roughly \(2\times\) per layer for ReLU
networks and is empirically a poor choice for from-scratch ResNet
training.

\paragraph{Implementation.} The
\code{\_init\_weights\_xavier} method of \code{AvancedResNet50TSM} was
renamed to \code{\_init\_weights} and rewritten to walk every submodule:

\begin{lstlisting}[style=py]
def _init_weights(self) -> None:
    for module in self.modules():
        if isinstance(module, nn.Conv2d):
            nn.init.kaiming_normal_(module.weight,
                mode="fan_out", nonlinearity="relu")
            if module.bias is not None:
                nn.init.zeros_(module.bias)
        elif isinstance(module, nn.BatchNorm2d):
            if module.weight is not None:
                nn.init.ones_(module.weight)
            if module.bias is not None:
                nn.init.zeros_(module.bias)
    nn.init.normal_(self.classifier.weight, mean=0.0, std=0.01)
    if self.classifier.bias is not None:
        nn.init.zeros_(self.classifier.bias)
\end{lstlisting}

The classifier head is initialised with a small-Gaussian
(\(\mathcal{N}(0, 0.01^2)\)) following the original ResNet and TSM
recipes; BN affine parameters are set to \(\gamma=1, \beta=0\).

\subsection{Automatic Mixed Precision (AMP)}
\label{sec:phase1:amp}

\paragraph{Logic.} On Ampere, Hopper and newer GPUs, running the forward
pass in FP16 (or BF16) and only the master copy of the optimizer state in
FP32 yields \(1.5\)--\(2\times\) throughput at no measurable accuracy
cost when paired with PyTorch's gradient scaler.

\paragraph{Implementation.} A new boolean \cfgkey{training.amp}
(default \code{false}) gates a \code{torch.amp.GradScaler} that is
created in \file{train.py} and threaded through the engine helpers.
\code{train\_one\_epoch} and \code{evaluate\_epoch} both received two new
keyword arguments, \code{scaler} and \code{amp\_dtype}, and wrap their
forward and backward passes in
\code{torch.amp.autocast(device\_type="cuda", dtype=amp\_dtype)}. When
\code{scaler} is \code{None} (the default), the code path is
byte-for-byte identical to the legacy full-precision implementation.

\begin{lstlisting}[style=py]
amp_enabled = bool(cfg.training.get("amp", False)) \
              and device.type == "cuda"
scaler = torch.amp.GradScaler("cuda", enabled=True) \
         if amp_enabled else None

# ...inside train_one_epoch:
with torch.amp.autocast(device_type="cuda",
                        enabled=use_amp,
                        dtype=amp_dtype):
    logits = model(videos)
    loss = loss_fn(logits, target)
if use_amp:
    scaler.scale(loss).backward()
    scaler.step(optimizer)
    scaler.update()
else:
    loss.backward()
    optimizer.step()
\end{lstlisting}

\subsection{Persistent optimizer / scheduler / scaler state}
\label{sec:phase1:resume}

\paragraph{Logic.} A long training run inevitably gets restarted
(scheduled maintenance, a transient OOM, an accidental \code{Ctrl-C}). If
only the model weights are persisted, the optimizer's running moments,
the cosine LR schedule and the GradScaler all reset, which causes a
visible accuracy dip for several epochs after each restart and makes
``resume from epoch~$k$'' a much weaker recovery than ``never
interrupted''.

\paragraph{Implementation.} On every \emph{best-checkpoint} save, the
trainer now stores \code{optimizer.state\_dict()},
\code{cosine\_scheduler.state\_dict()} (when present) and
\code{scaler.state\_dict()} (when AMP is on) inside the existing
\code{extra} dict of the checkpoint payload. The schema version is
unchanged. On resume, those state dicts are loaded inside try/except
blocks so a checkpoint produced by an older code path still resumes
cleanly (the cosine scheduler then falls back to its legacy
fast-forward).

\begin{lstlisting}[style=py]
ckpt_extra: dict[str, Any] = {
    "val_top1": ckpt_stats.top1,
    "val_top5": ckpt_stats.top5,
    "val_loss": ckpt_stats.loss,
    "epoch": epoch + 1,
    "optimizer_state_dict": optimizer.state_dict(),
    "trained_class_indices": trained_class_indices,
    "checkpoint_kind": ckpt_kind,
}
if cosine_scheduler is not None:
    ckpt_extra["scheduler_state_dict"] = cosine_scheduler.state_dict()
if scaler is not None:
    ckpt_extra["scaler_state_dict"] = scaler.state_dict()
\end{lstlisting}

\subsection{Trained-class indices and untrained-logit masking}
\label{sec:phase1:masking}

\paragraph{Logic.} The provided dataset has 33 unique class indices in
\([0, 32]\) but class~27 has no training samples. The classifier head is
nonetheless 33-wide (its weight tensor was initialised at random for
class~27) and at inference time argmax can therefore land on class~27 by
accident, producing systematically wrong predictions for every test
clip on which the random class~27 logit happens to be the largest.

\paragraph{Implementation.} The trainer derives the set of class indices
that actually have at least one training sample and writes it into the
checkpoint:

\begin{lstlisting}[style=py]
trained_class_indices: list[int] = sorted(
    {int(label) for _, label in train_samples})
\end{lstlisting}

The submission pipeline reads this list back and builds an additive mask
that pushes every never-trained class to \(-\infty\) before
\code{argmax}:

\begin{lstlisting}[style=py]
def _build_untrained_mask(trained_class_indices, num_classes, device):
    if not trained_class_indices:
        return None
    trained = {int(i) for i in trained_class_indices
               if 0 <= int(i) < num_classes}
    if len(trained) >= num_classes:
        return None
    mask = torch.zeros(num_classes, dtype=torch.float32, device=device)
    untrained = [c for c in range(num_classes) if c not in trained]
    mask[untrained] = -math.inf
    return mask
\end{lstlisting}

Older checkpoints that pre-date Phase~1 do not carry this list; in that
case the mask is \code{None} and the legacy unmasked argmax path is used,
preserving backward compatibility.

% =====================================================================
\section{Architectural and Training Enhancements}
\label{sec:phase2}

Phase~2 adds four orthogonal opt-in features that compose cleanly with
Phase~1. None of them change the behaviour of the legacy code paths when
left at their defaults; together, they are bundled in a new experiment
config \file{configs/experiment/track\_a\_phase2.yaml} that also flips on
AdamW with a stronger weight decay.

\subsection{Stochastic Depth (DropPath)}
\label{sec:phase2:droppath}

\paragraph{Logic.} Stochastic Depth \cite{Huang2016SD} is a stochastic
regularisation technique that, during training, drops the residual branch
of each ResNet bottleneck block with probability \(p_\ell\) and rescales
the kept contributions by \(1/(1-p_\ell)\) so that the expectation over
the random mask matches the no-drop forward pass:
\[
y_\ell \;=\; x_\ell \;+\;
\frac{b_\ell}{1-p_\ell}\,F_\ell(x_\ell), \qquad
b_\ell \sim \text{Bernoulli}(1-p_\ell).
\]
At inference time, all blocks are kept (\(b_\ell = 1\) deterministically)
so no rescaling is needed. The standard recipe uses a linear schedule
that drops deeper blocks more aggressively:
\[
p_\ell \;=\; p_{\max}\,\frac{\ell - 1}{L - 1},
\]
where \(L\) is the total number of bottleneck blocks. For ResNet-50
\(L=16\) and a typical \(p_{\max}\) is in \([0.05, 0.2]\).

\paragraph{Implementation.} A parameter-free \code{DropPath} module
implements the per-sample bernoulli mask. A helper
\code{\_inject\_drop\_path} walks the four ResNet stages, attaches a
\code{DropPath} instance to each block (with the linearly-scaled
probability), and rebinds the block's \code{forward} to a patched
implementation that calls \code{self.drop\_path(out)} on the residual
branch immediately before the residual addition:

\begin{lstlisting}[style=py]
def _patched_bottleneck_forward(self, x: torch.Tensor) -> torch.Tensor:
    identity = x
    out = self.relu(self.bn1(self.conv1(x)))
    out = self.relu(self.bn2(self.conv2(out)))
    out = self.bn3(self.conv3(out))
    if self.downsample is not None:
        identity = self.downsample(x)
    out = self.drop_path(out)         # <-- new
    return self.relu(out + identity)
\end{lstlisting}

\code{DropPath} adds no parameters and no buffers, so the model's
\code{state\_dict} layout is unchanged regardless of
\cfgkey{model.drop\_path\_rate}. This is asserted by a unit test that
compares state-dict key sets with and without DropPath enabled.

\subsection{Attention-pool head}
\label{sec:phase2:attn}

\paragraph{Logic.} The legacy temporal head averages the per-frame
features:
\[
z = \frac{1}{T}\sum_{t=1}^{T} \phi(x_t), \qquad
\widehat{y} = W z + b.
\]
A 1-query multi-head attention pool is a strict generalisation: when the
attention weights are uniform it reproduces the mean exactly, but it can
also concentrate on the most informative frame(s). Concretely, we
introduce a learnable query token \(q \in \mathbb{R}^{1\times D}\) and
compute
\[
z \;=\; \mathrm{LayerNorm}\!\left(
        \mathrm{MultiHeadAttn}(q,\, K{=}V{=}\Phi)\right),
\]
where \(\Phi \in \mathbb{R}^{T\times D}\) is the stack of frame
features. With \(T=4\) and \(D=2048\), this adds about
\(D + 4D^2 \approx 16{,}780{,}288\) parameters, which is small compared
to the \(\approx 25\)M parameter ResNet-50 trunk.

\paragraph{Implementation.} A new \code{AttentionPool} module wraps a
single \code{nn.MultiheadAttention} layer plus a final
\code{nn.LayerNorm}; the query token is initialised with a truncated
normal of std~0.02. The supervised model exposes a new
\cfgkey{model.head} switch (\code{"mean"} or \code{"attn"}) and routes
features accordingly:

\begin{lstlisting}[style=py]
sequence_features = frame_features.view(batch_size, num_frames, -1)
if self.attn_pool is not None:
    consensus_features = self.attn_pool(sequence_features)
else:
    consensus_features = sequence_features.mean(dim=1)
\end{lstlisting}

When \cfgkey{model.head}~$=$~\code{"mean"} (the default) the
\code{attn\_pool} attribute is \code{None} and no extra parameters are
added.

\subsection{Exponential Moving Average (EMA) of model weights}
\label{sec:phase2:ema}

\paragraph{Logic.} Polyak averaging \cite{PolyakSIAM1992} maintains an
exponentially-weighted average of the model's parameters,
\[
\theta_{\mathrm{EMA}}^{(t+1)} \;=\;
\alpha\,\theta_{\mathrm{EMA}}^{(t)} \;+\;
(1-\alpha)\,\theta^{(t+1)}, \qquad \alpha \in [0.99, 0.9999],
\]
which empirically gives smoother loss curves and a small but consistent
top-1 boost (\(\approx 0.3{-}1.0\) points) on most modern image and
video classifiers.

\paragraph{Implementation.} We use PyTorch's
\code{torch.optim.swa\_utils.AveragedModel} with a custom averaging
function and \code{use\_buffers=True} so BatchNorm running statistics are
also tracked:

\begin{lstlisting}[style=py]
def _ema_avg_fn(avg_param, model_param, _num_averaged):
    return ema_decay * avg_param + (1.0 - ema_decay) * model_param

ema_model = torch.optim.swa_utils.AveragedModel(
    model, avg_fn=_ema_avg_fn, use_buffers=True).to(device)
\end{lstlisting}

\code{train\_one\_epoch} accepts a new \code{ema\_model} keyword and
calls \code{ema\_model.update\_parameters(model)} after every optimizer
step. At the end of each epoch the trainer evaluates \emph{both} the
live and the EMA model and saves the better of the two:

\begin{lstlisting}[style=py]
if ema_stats is not None and ema_stats.top1 > val_stats.top1:
    ckpt_stats, ckpt_module, ckpt_kind = ema_stats, ema_model, "ema"
else:
    ckpt_stats, ckpt_module, ckpt_kind = val_stats, model, "live"

# ...later, before save_checkpoint:
module_to_save = (ckpt_module.module
                  if isinstance(ckpt_module,
                                torch.optim.swa_utils.AveragedModel)
                  else ckpt_module)
\end{lstlisting}

This selection rule is what lets downstream consumers
(\file{submit.py}, \file{evaluate.py}) stay completely oblivious to the
EMA toggle: the saved \code{model\_state\_dict} is always loaded with
\code{build\_model(...).load\_state\_dict(...)}, regardless of which
copy won. The string \code{"live"} or \code{"ema"} is recorded in
\code{extra["checkpoint\_kind"]} for traceability.

\subsection{Test-Time Augmentation with auto class-pair remap}
\label{sec:phase2:tta}

\paragraph{Logic.} Horizontal flip is a free \emph{evaluation-time}
augmentation \emph{provided} that the predicted class is invariant to
the flip. The Something-Something~V2 subset used in this competition
contains a small number of direction-sensitive classes, in particular:
\begin{itemize}[leftmargin=2em]
\item class~018: ``Pulling something from left to right''
\item class~019: ``Pulling something from right to left''
\end{itemize}
For these, a flipped clip should predict the \emph{paired} class. A
naive flip-and-average TTA actually hurts accuracy on these two classes
because it averages two confident but \emph{disagreeing} predictions.

\paragraph{Implementation.} We average softmax probabilities across the
two views (original and flipped) and apply a learned permutation
\(\pi: \{0,\ldots,K-1\}\to\{0,\ldots,K-1\}\) to the flipped softmax so
that paired classes are mapped back into agreement. The permutation is
\emph{auto-derived from the class folder names}, making it robust to any
future change in the class set:

\begin{lstlisting}[style=py]
LR = "from_left_to_right"
RL = "from_right_to_left"
for idx, stem in stem_by_idx.items():
    if LR in stem:
        mirror_stem = stem.replace(LR, RL)
        for other_idx, other_stem in stem_by_idx.items():
            if other_stem == mirror_stem:
                perm[idx] = other_idx
                perm[other_idx] = idx
                break
\end{lstlisting}

(The numeric prefix is stripped before matching so
\code{018\_..\_left\_to\_right} pairs with \code{019\_..\_right\_to\_left}
without a code-level hard-code.) The TTA loop in \file{submit.py} then
becomes:

\begin{lstlisting}[style=py]
logits = _logits_for_batch(model, video_batch, untrained_mask)
probs_total = torch.softmax(logits, dim=1)
n_views = 1
if tta_flip:
    flipped = torch.flip(video_batch, dims=[-1])
    flipped_probs = torch.softmax(
        _logits_for_batch(model, flipped, untrained_mask), dim=1)
    if flip_perm is not None:
        flipped_probs = flipped_probs.index_select(dim=1, index=flip_perm)
    probs_total = probs_total + flipped_probs
    n_views += 1
predictions.extend(int(p) for p in
    (probs_total / n_views).argmax(dim=1).cpu().tolist())
\end{lstlisting}

The configuration is opt-in (\cfgkey{training.tta}, \cfgkey{training.tta\_flip},
both \code{false} by default) and read at inference time only; training
is unaffected.

\subsection{AdamW and the bundled experiment configuration}
\label{sec:phase2:adamw}

The trainer already supported AdamW; Phase~2 simply makes it the
recommended optimizer for the new experiment by setting
\cfgkey{training.optimizer: adamw} and \cfgkey{training.weight\_decay:
0.05} in \file{configs/experiment/track\_a\_phase2.yaml}. Decoupling the
weight-decay term from the gradient (which is what AdamW does, unlike
plain Adam) is the right thing to do whenever a non-zero weight decay is
used; it is also the optimizer of choice for ConvNeXt, ViT and modern
ResNet recipes.

The \file{track\_a\_phase2.yaml} bundle composes all of Phase~1 and
Phase~2:

\begin{lstlisting}[style=yaml]
defaults:
  - override /model: avanced_resnet50_tsm
  - override /augment: randaugment_t

model:
  drop_path_rate: 0.1
  head: attn
  head_num_heads: 4

dataset:
  num_frames: 4
  use_official_val: true

training:
  optimizer: adamw
  weight_decay: 0.05
  amp: true
  ema_enabled: true
  ema_decay: 0.999
  tta: true
  tta_flip: true
  scheduler_cosine: true
  warmup_epochs: 10
  label_smoothing: 0.1
\end{lstlisting}

\section{Self-Supervised Pretraining (V-JEPA)}
\label{sec:phase3}

The single highest-impact change available within Track~A is to give the
trunk something useful to do with the unlabeled clips before the
supervised loss starts back-propagating. Track~A permits self-supervised
pretraining on the \emph{provided} unlabeled images (the labels are
simply not consumed and the trunk starts from random initialisation), so
we pretrain on the union of the \file{train/}, \file{val/} and
\file{test/} folders.

\subsection{Why the DINO-on-still-frames prototype was deprecated}
\label{sec:phase3:dino-deprecated}

Our first SSL attempt was a DINO-v1 \cite{Caron2021DINO} pipeline trained
on still frames sampled independently from the provided videos (file
\file{src/smth2smth/shared/data/still\_frames\_dataset.py} and
\file{src/smth2smth/shared/models/dino\_ssl.py}). In practice it produced
\emph{inconsistent} downstream gains --- sometimes a small improvement,
sometimes a clear regression versus the random-init baseline. Two
structural issues explain why:

\begin{enumerate}[leftmargin=*]
\item \textbf{DINO is invariance-trained.} The objective is to assign
   the same prototype to two augmented views of the \emph{same image},
   which actively pressures the trunk to become invariant to the small
   per-pixel changes that distinguish adjacent video frames --- exactly
   the signal an action-anticipation model needs to keep.
\item \textbf{The TSM modules are never exercised.} DINO operates on
   isolated frames, so the channel-shift modules wrapped around every
   bottleneck's \code{conv1} (which carry $\approx$ 12.5\% of the
   input channels temporally) see no temporal context during
   pretraining. Roughly one in eight conv1 input channels was therefore
   left at its Kaiming He initial value when supervised training
   started.
\end{enumerate}

The DINO files remain in the repository as a historical record, but the
recommended SSL pipeline is now the V-JEPA-style one described in the
rest of this section.

\subsection{V-JEPA-style clip-level pretraining}
\label{sec:phase3:vjepa-overview}

V-JEPA (Joint-Embedding Predictive Architecture for Video)
\cite{Bardes2024VJEPA, Assran2025VJEPA2} replaces the
``invariance to augmented views'' objective by ``predict the masked
parts of a clip in feature space''. Concretely, the student encodes a
\emph{masked} clip while the EMA-updated teacher encodes the unmasked
clip; a small predictor maps the student's representation to the
teacher's space, and an L1 loss is applied only at the masked positions.
This is the right inductive bias for our problem: the student is
forced to use temporal context to fill in masked features, which makes
the TSM channel-shift modules immediately useful.

We adapt the recipe to our TSM-equipped ResNet-50 trunk with $T=4$
frames. Because there are only 4 frames per clip, full V-JEPA 3D
multi-block tube masking in patch-token space is not appropriate; we
use \emph{frame-level} masking instead, which mirrors the V-JEPA
objective at the granularity that maps naturally onto our CNN trunk.

\subsection{Clip dataset for SSL}
\label{sec:phase3:dataset}

A new file \file{src/smth2smth/shared/data/clip\_ssl\_dataset.py}
exposes two utilities:

\begin{itemize}[leftmargin=2em]
\item \code{collect\_all\_video\_dirs(roots)} walks the provided
      directories and returns every video folder it finds. It handles
      both layouts in use: class-bucketed
      (\file{train/<class>/<video>/<frame>.jpg}) and flat
      (\file{test/<video>/<frame>.jpg}). Duplicates across roots are
      removed.
\item \code{ClipSSLDataset} yields one \((T, C, H, W)\) clip per video.
      Frame indices are picked by the same \code{pick\_frame\_indices}
      helper used by supervised training so the SSL run sees exactly
      the same temporal distribution as the supervised trainer. The
      per-frame transform is the project's standard
      \code{build\_transforms(is\_training=True)}, which includes
      RandAugment-T spatial augmentations \emph{without} ImageNet
      normalisation (Track~A forbids pretrained weights).
\end{itemize}

\paragraph{On the question of V-JEPA's data augmentations.} The V-JEPA
paper combines random resized crop, horizontal flip, colour jitter,
gray-scale and Gaussian blur in a specific stack. We \emph{do not}
reproduce that augmentation stack verbatim: we reuse the project's
existing RandAugment-T pipeline so that the SSL trunk sees exactly the
same input distribution as the downstream supervised trainer.
Reproducing V-JEPA's stack would put a domain shift between pretraining
and fine-tuning, which is the opposite of what we want. The
\emph{augmentation} contribution of V-JEPA that we \emph{do} reproduce
is the masking itself: \(n_{\mathrm{masked}}\) frames per clip are
replaced by zeros so the student must reconstruct their features from
the unmasked context through the TSM channels.

\subsection{Model: trunk + predictor + EMA teacher}
\label{sec:phase3:model}

\file{src/smth2smth/shared/models/vjepa\_ssl.py} contains three pieces.

\paragraph{Trunk.} \code{VJepaTrunk} is the \emph{same}
TSM-wrapped ResNet-50 as the supervised model's \code{backbone}: it
calls \code{\_make\_temporal\_shift\_resnet} with \code{n\_segment} =
\(T\), so every bottleneck block's \code{conv1} is replaced by
\code{nn.Sequential(TemporalShift, conv1)}. The state-dict layout
therefore matches \code{AvancedResNet50TSM.backbone.state\_dict()}
byte-for-byte, which lets the supervised trainer pick up the SSL
checkpoint with a single \code{load\_state\_dict(..., strict=False)}.

\paragraph{Predictor.} A small per-frame MLP
(\code{Linear} $\to$ \code{GELU} $\to$ \code{Linear}, hidden dim 2048)
maps the student's per-frame features into the teacher's space. The
predictor exists so the encoder is not biased toward an identity
student-to-teacher mapping on masked positions (Bardes et al., Sec.~3).

\paragraph{Composite.} \code{VJepaModel} composes the trunk and the
predictor. Two copies are instantiated at training time --- a student
and a teacher --- with identical architecture; the teacher's weights
are an EMA of the student's and are never updated by gradient descent.

\subsection{Frame-level masking and the V-JEPA L1 loss}
\label{sec:phase3:loss}

For each clip we sample \(n_{\mathrm{masked}} \sim
\mathrm{Uniform}[1, T-1]\) and zero out the corresponding frames. The
student trunk encodes the masked clip; the teacher trunk encodes the
unmasked clip with no gradient. The masked-position L1 loss is
\[
\mathcal{L} \;=\;
\frac{1}{|\mathcal{M}| \cdot D}\sum_{(b,t) \in \mathcal{M}}
\left\| \hat{f}_{b,t} - f^{\mathrm{tch}}_{b,t} \right\|_1,
\]
where \(\hat{f}_{b,t} = \mathrm{Predictor}(\mathrm{Student}(\tilde
x)_{b,t})\) is the predicted feature for clip \(b\) at frame \(t\),
\(f^{\mathrm{tch}}_{b,t} = \mathrm{Teacher}(x)_{b,t}\) is the teacher's
target, \(\mathcal{M}\) is the set of masked \((b, t)\) positions and
\(D=2048\). Masked positions are the only ones that contribute to the
loss; unmasked frames see no learning signal, which keeps the gradient
focused on \emph{predicting what is missing}.

\begin{lstlisting}[style=py]
def vjepa_feature_loss(predicted, teacher, mask):
    if not mask.any():
        return predicted.sum() * 0.0  # zero with valid grad_fn
    diff = (predicted - teacher).abs()        # (B, T, D)
    sel = mask.to(diff.dtype).unsqueeze(-1)   # (B, T, 1)
    weighted = diff * sel
    n_elems = mask.sum().clamp_min(1).to(diff.dtype) * diff.shape[-1]
    return weighted.sum() / n_elems
\end{lstlisting}

\subsection{Teacher EMA momentum schedule}

The teacher's parameters are updated by
\(\theta_t \leftarrow m\,\theta_t + (1-m)\,\theta_s\) after every
optimizer step, with \(m\) following a cosine schedule from
\(m_0 = 0.996\) at the start to \(m_1 = 1.0\) at the end of training.

\subsection{Trunk-only checkpoint format}

After every epoch we save a small ($\approx$ 95~MiB) checkpoint that
contains \emph{only} the trunk parameters, with both the \code{trunk.}
and the inner \code{backbone.} prefix stripped. The result is a flat
dictionary keyed exactly like a standard
\code{torchvision.models.resnet50} state-dict, except that the
bottleneck \code{conv1} keys already carry the TSM-wrapping index (e.g.
\code{layer1.0.conv1.1.weight} instead of
\code{layer1.0.conv1.weight}):

\begin{lstlisting}[style=py]
trunk_state_dict = {}
for k, v in student.state_dict().items():
    if k.startswith("trunk.backbone."):
        trunk_state_dict[k.removeprefix("trunk.backbone.")] = v
torch.save({"trunk_state_dict": trunk_state_dict,
            "epoch": epoch + 1}, out_path)
\end{lstlisting}

\subsection{The \cfgkey{model.init\_from} hook in \file{train.py}}
\label{sec:phase3:initfrom}

The supervised trainer received a single new opt-in:

\begin{lstlisting}[style=py]
init_from = cfg.model.get("init_from") if hasattr(cfg.model, "get") \
            else None
if init_from:
    payload = torch.load(Path(str(init_from)).resolve(),
                         map_location=device, weights_only=False)
    trunk_state = payload.get("trunk_state_dict")
    prefixed = _ssl_trunk_to_supervised_keys(trunk_state)
    missing, unexpected = model.load_state_dict(prefixed, strict=False)
\end{lstlisting}

The helper \code{\_ssl\_trunk\_to\_supervised\_keys} handles both
flavours of SSL checkpoint that the repository can produce:

\begin{itemize}[leftmargin=2em]
\item Legacy DINO trunks whose keys come from a plain
   \code{torchvision} ResNet-50 (e.g. \code{layer1.0.conv1.weight}).
   The helper renames them to the TSM-wrapped form
   (\code{layer1.0.conv1.1.weight}) and prepends \code{backbone.}.
\item V-JEPA trunks whose keys are already TSM-wrapped
   (\code{layer1.0.conv1.1.weight}). The rename regex is a no-op and
   the helper only prepends \code{backbone.}.
\end{itemize}

\begin{lstlisting}[style=py]
block_conv1_re = re.compile(
    r"^(layer[1-4]\.\d+\.conv1)\.(weight|bias)$")
out: dict[str, torch.Tensor] = {}
for k, v in trunk_state.items():
    match = block_conv1_re.match(k)
    if match is not None:
        new_k = f"{match.group(1)}.1.{match.group(2)}"
    else:
        new_k = k
    out[f"backbone.{new_k}"] = v
return out
\end{lstlisting}

A V-JEPA trunk loads with \textbf{318 backbone tensors covered, zero
unexpected keys, and only the classifier (and, when enabled, the
attention pool) left at their fresh-init values} --- the exact behaviour
the test suite pins (\file{tests/pipelines/test\_init\_from\_vjepa.py}).

% =====================================================================
\section{Data Augmentation, Class Balance, and TTA Improvements}
\label{sec:tracka-phase4}

Phase~4 attacks the failure modes that remain once a good trunk is in
place. Three of them stem directly from the small, imbalanced nature of
the provided training set; the fourth is a free accuracy boost at
inference time. All four are opt-in via the configuration and are
bundled together in
\file{configs/experiment/track\_a\_vjepa\_finetune.yaml}.

\subsection{Class imbalance: balanced sampling \emph{and} balanced loss}
\label{sec:tracka-phase4:classbalance}

\paragraph{Logic.} The training set is heavily skewed: class counts
range from \(\approx\) 162 to \(\approx\) 3170 samples, a 20\(\times\)
ratio. A vanilla cross-entropy loss with uniform sampling effectively
\emph{down-weights} rare classes twice over (they are seen less often,
\emph{and} each sample contributes equally to the loss), and the
resulting classifier tends to predict the head classes for most test
clips. We fix this on two independent axes.

\paragraph{Balanced sampler.} We replace the uniform shuffle of the
\code{DataLoader} by a \code{WeightedRandomSampler} whose per-sample
weight is \(w_i = 1 / \sqrt{n_{c(i)}}\), where \(c(i)\) is the class
of sample \(i\) and \(n_c\) is the class count. The square root is a
robust compromise between fully-balanced sampling (which heavily
oversamples the 162-sample classes) and uniform sampling (which
ignores them). The policy is selected by
\cfgkey{training.class\_balance\_sampler} \(\in
\{\text{none, inverse, sqrt\_inverse}\}\).

\paragraph{Balanced loss.} We additionally weight each class in the
cross-entropy by Cui et al.'s effective-number weights
\cite{Cui2019CB}: for class \(c\) with \(n_c\) samples,
\[
w_c \;\propto\; \frac{1-\beta}{1-\beta^{n_c}}, \qquad \beta \in (0, 1),
\]
which interpolates between uniform weighting (\(\beta \to 0\)) and
inverse-frequency weighting (\(\beta \to 1\)). We use the canonical
\(\beta = 0.999\). Classes with zero samples (e.g. class~27, see
Section~\ref{sec:phase1:masking}) receive \(w_c = 0\), so they
contribute neither to the sampling nor to the loss. The same weight
vector is also routed through the soft-target cross-entropy used by
label smoothing and FrameMixUp so the augmentations remain class-balanced.

\begin{lstlisting}[style=py]
# Per-class effective-number weights (Cui et al. 2019).
def _class_factors(n_per_class, policy, beta=0.999):
    if policy == "cb":
        effective = [(1.0 - beta**max(n, 1)) / (1.0 - beta)
                     for n in n_per_class]
        return [1.0 / e if e > 0 else 0.0 for e in effective]
    # ... 'inverse', 'sqrt_inverse' branches ...
\end{lstlisting}

The two axes are configured independently
(\cfgkey{training.class\_balance\_sampler},
\cfgkey{training.class\_balance\_loss},
\cfgkey{training.class\_balance\_beta}) so the user can run any
combination of \(\{\text{sampler off, sampler sqrt}\} \times
\{\text{loss off, loss CB}\}\) and compare. The new
\code{class\_balance.py} module is fully unit-tested
(\file{tests/shared/utils/test\_class\_balance.py}).

\subsection{Motion-aware time-reversal augmentation}
\label{sec:tracka-phase4:timereversal}

\paragraph{Logic.} For verbs whose semantics are intrinsically
direction-sensitive, the time-reversed clip belongs to a \emph{paired}
class --- the reverse of opening is closing, the reverse of pushing
left-to-right is pushing right-to-left, etc. Time-reversal is therefore
a label-preserving augmentation \emph{only} when we also remap the
label to the paired class. Without the remap, a reversed clip injects
flat-out wrong supervision.

\paragraph{Implementation.} A new module
\file{src/smth2smth/shared/data/time\_reversal.py} hosts a small
hand-curated table of reversible class pairs:
\begin{quote}
\(\text{Opening} \leftrightarrow \text{Closing}\),
\(\text{Moving up} \leftrightarrow \text{Moving down}\),
\(\text{Pushing left-to-right} \leftrightarrow
\text{Pushing right-to-left}\),
\(\text{Tilting toward camera} \leftrightarrow
\text{Tilting away from camera}\),
plus a symmetric set (e.g. \emph{Holding something}) where the reverse
is the same class.
\end{quote}
\code{build\_time\_reversal\_table} matches the pairs against the
actual class folder names on disk and returns two tensors:
\code{perm[c]} = the class index of the time-reversed clip
(\code{perm[c] = -1} if \(c\) is not reversible), and
\code{allow\_mask[c] = True} iff class \(c\) is reversible. Both are
threaded into \code{VideoFrameDataset} via three new keyword
arguments:

\begin{lstlisting}[style=py]
class VideoFrameDataset(Dataset):
    def __init__(self, ..., *,
                 time_reversal_prob: float = 0.0,
                 time_reversal_perm: torch.Tensor | None = None,
                 time_reversal_allow_mask: torch.Tensor | None = None):
        ...
    def __getitem__(self, index):
        ...
        if self._should_reverse(label):
            indices = list(reversed(indices))
            label = int(self.time_reversal_perm[int(label)].item())
        ...
\end{lstlisting}

The augmentation is gated by \cfgkey{dataset.time\_reversal\_prob}
(default \(0.0\)). When set to e.g. \(0.3\), about 30\% of clips drawn
from reversible classes are reversed and re-labelled at every epoch,
which roughly doubles the effective training set for the affected
classes without any synthetic data.

\subsection{Multi-scale TTA and logit adjustment}
\label{sec:tracka-phase4:tta}

The Phase~2 TTA (Section~\ref{sec:phase2:tta}) already averages
horizontal-flip predictions with an auto-derived left/right class
permutation. Phase~4 extends inference with two further tricks.

\paragraph{Multi-scale TTA.} At inference time we resample the input
clip to several spatial scales (e.g. \(\{0.875, 1.0, 1.125\} \times
H \times W\)) using bilinear interpolation, run the model on each
scale (optionally with a horizontal flip), softmax each set of
logits, and average all softmaxes. This is a standard test-time
ensembling trick for CNN classifiers; it is gated by
\cfgkey{training.tta\_scales} (a list, default \([1.0]\)).

\begin{lstlisting}[style=py]
def _rescale_video(video_batch, scale):
    if abs(scale - 1.0) < 1e-6:
        return video_batch
    b, t, c, h, w = video_batch.shape
    new_h, new_w = int(round(h * scale)), int(round(w * scale))
    return F.interpolate(video_batch.reshape(b * t, c, h, w),
                         size=(new_h, new_w),
                         mode="bilinear",
                         align_corners=False).view(b, t, c, new_h, new_w)
\end{lstlisting}

\paragraph{Logit adjustment.} Even with a class-balanced sampler and
loss, the predictor at test time still inherits a small bias toward
head classes whenever the training data and the test data have
different (but unknown) class distributions. Menon et al.
\cite{Menon2021LA} show that subtracting \(\tau \cdot \log \pi_c\) from
class \(c\)'s logit at \emph{inference time}, where \(\pi_c\) is the
training-set prior and \(\tau \in [0, 1]\) is a small constant, is the
Bayes-optimal post-hoc correction under the standard balanced-test
assumption. We compute \(\pi_c\) directly from the training-folder
counts and subtract \(\tau \log \pi_c\) from the logits right after the
forward pass, before either the softmax or the multi-scale averaging:

\begin{lstlisting}[style=py]
def _build_logit_adjustment(train_dir, num_classes, tau, device):
    counts = class_counts(...)
    total = float(sum(counts))
    pi = torch.tensor([c / total if c > 0 else 1e-12 for c in counts],
                      device=device)
    return tau * torch.log(pi)        # subtracted from logits

# inside _logits_for_batch:
logits = model(video_batch)
if logit_adjust is not None:
    logits = logits - logit_adjust
\end{lstlisting}

The strength \(\tau\) is configured via
\cfgkey{training.tta\_logit\_adjust} (default \(0.0\) = off; we use
\(\tau = 1.0\) in the bundled experiment).

\subsection{The bundled experiment configuration}

The four Phase-4 features compose cleanly with everything in Phases~1
through~3. A new experiment file
\file{configs/experiment/track\_a\_vjepa\_finetune.yaml} bundles them on
top of a V-JEPA warm-start:

\begin{lstlisting}[style=yaml]
defaults:
  - override /model: avanced_resnet50_tsm
  - override /augment: randaugment_t

model:
  drop_path_rate: 0.1
  head: attn
  head_num_heads: 4
  init_from: null      # set on the command line

dataset:
  num_frames: 4
  use_official_val: true
  time_reversal_prob: 0.3

training:
  optimizer: adamw
  weight_decay: 0.05
  amp: true
  ema_enabled: true
  ema_decay: 0.999
  tta: true
  tta_flip: true
  tta_scales: [0.875, 1.0, 1.125]
  tta_logit_adjust: 1.0
  class_balance_sampler: sqrt_inverse
  class_balance_loss: cb
  class_balance_beta: 0.999
  early_stopping_enabled: true
  early_stopping_patience: 20
\end{lstlisting}

% =====================================================================
\section{Configuration and reproducibility}
\label{sec:configs}

\subsection{New configuration knobs}

The following table summarises every new switch introduced by Phases~1
through~4. All defaults preserve the legacy behaviour.

\begin{center}\small
\begin{tabular}{@{}lll@{}}
\toprule
\textbf{Key} & \textbf{Default} & \textbf{Effect when on} \\
\midrule
\multicolumn{3}{l}{\emph{Phases 1--2}} \\
\cfgkey{dataset.use\_official\_val}  & \code{false} & validate on \file{data/val/} \\
\cfgkey{training.amp}                & \code{false} & FP16 autocast + GradScaler \\
\cfgkey{training.ema\_enabled}       & \code{false} & EMA of model weights \\
\cfgkey{training.ema\_decay}         & \code{0.999} & EMA momentum \\
\cfgkey{training.tta}                & \code{false} & enable TTA at submit time \\
\cfgkey{training.tta\_flip}          & \code{true}  & flip view (only used when \code{tta}) \\
\cfgkey{training.tta\_temporal\_shift}\!\! & \code{0} & extra temporal start-offsets (TTA) \\
\cfgkey{model.drop\_path\_rate}      & \code{0.0}   & Stochastic Depth maximum drop prob.\\
\cfgkey{model.head}                  & \code{mean}  & temporal head (\code{mean}/\code{attn}) \\
\cfgkey{model.head\_num\_heads}      & \code{4}     & MultiHead attention heads \\
\cfgkey{model.init\_from}            & \code{null}  & path to SSL trunk to warm-start \\
\midrule
\multicolumn{3}{l}{\emph{Phase 3 (V-JEPA SSL)}} \\
\cfgkey{pretrain.mask\_prob}         & \code{0.5}   & avg.\ frame mask fraction for V-JEPA \\
\cfgkey{pretrain.min\_mask}          & \code{1}     & min frames masked per clip \\
\cfgkey{pretrain.max\_mask}          & \code{2}     & max frames masked per clip \\
\cfgkey{pretrain.predictor\_hidden\_dim}\!\! & \code{2048} & MLP predictor hidden width \\
\cfgkey{pretrain.teacher\_momentum\_base}\!\! & \code{0.996} & cosine EMA momentum start \\
\cfgkey{pretrain.teacher\_momentum\_final}\!\! & \code{1.0} & cosine EMA momentum end \\
\midrule
\multicolumn{3}{l}{\emph{Phase 4}} \\
\cfgkey{training.class\_balance\_sampler}\!\! & \code{none} & sample weights: \code{none/inverse/sqrt\_inverse} \\
\cfgkey{training.class\_balance\_loss}\!\! & \code{none} & class-balanced loss policy (\code{none/cb}) \\
\cfgkey{training.class\_balance\_beta}\!\! & \code{0.999} & \(\beta\) in the CB effective-number formula \\
\cfgkey{dataset.time\_reversal\_prob}\!\! & \code{0.0} & prob.\ of time-reversing a reversible-class clip \\
\cfgkey{training.tta\_scales}\!\!     & \code{[1.0]} & list of spatial scales for multi-scale TTA \\
\cfgkey{training.tta\_logit\_adjust}\!\! & \code{0.0} & \(\tau\) in Menon et al.\ logit adjustment \\
\bottomrule
\end{tabular}
\end{center}

\subsection{New experiment configurations}

\begin{description}[leftmargin=1.5em, style=nextline]
\item[\file{configs/experiment/track\_a\_phase2.yaml}]
   Bundles every Phase-1 and Phase-2 opt-in: AdamW + WD~0.05, AMP,
   DropPath~0.1, attention head, EMA~0.999, TTA at submit time, official
   validation split, cosine schedule with 10-epoch warmup, label
   smoothing~0.1, RandAugment-T, FrameMixUp.

\item[\file{configs/experiment/track\_a\_phase2\_balanced.yaml}]
   Phase~2 \emph{plus} all of Phase~4 (class-balanced sampler + CB loss,
   time-reversal augmentation, multi-scale TTA, logit adjustment) but
   \emph{without} SSL warm-start. Useful as an ablation baseline.

\item[\file{configs/experiment/track\_a\_vjepa\_pretrain.yaml}]
   Routes the run through \code{pretrain\_vjepa.py}: 50 epochs of
   V-JEPA-style clip mask-denoising over the union of train+val+test
   videos, 1--2 frames masked per clip, AdamW + WD~0.04, AMP on.

\item[\file{configs/experiment/track\_a\_vjepa\_finetune.yaml}]
   Same as \file{track\_a\_phase2\_balanced.yaml} but expects
   \cfgkey{model.init\_from} to be set on the command line to the
   V-JEPA trunk produced by the previous experiment; reduces the
   number of supervised epochs (warm-started runs converge faster) and
   halves the learning rate.

\item[\file{configs/experiment/track\_a\_ssl\_pretrain.yaml},
      \file{...\_ssl\_finetune.yaml}]
   The legacy DINO-on-still-frames experiments are retained for
   reference. They are superseded by the V-JEPA variants above and are
   not recommended for new runs (see
   Section~\ref{sec:phase3:dino-deprecated}).
\end{description}

\subsection{Backwards compatibility}

Every legacy entry point continues to work without modification:

\begin{itemize}[leftmargin=2em]
\item Phase-1 checkpoints (no \code{trained\_class\_indices}, no
      \code{optimizer\_state\_dict} in \code{extra}) load and resume via
      the legacy fast-forward fallback.
\item Pre-Phase-1 checkpoints (no Phase-2 opt-ins enabled) load with
      identical state-dict keys; this is enforced by a unit test that
      compares state-dict key sets between defaults and Phase-2
      opt-ins.
\item Submission CSVs produced before Phase~1 remain valid; the new
      mask is only applied if the checkpoint actually carries the
      trained-class list.
\end{itemize}

% =====================================================================
\section{Test coverage}
\label{sec:tests}

The pre-existing suite contained 120 tests, all of which still pass.
Forty-one new tests were added across the four phases, bringing the
total to \textbf{161 passing tests}:

\begin{description}[leftmargin=1.5em, style=nextline]
\item[\file{tests/shared/models/test\_phase2\_optins.py} (8 tests)]
   Phase~2: \code{DropPath} eval-time identity and unbiased
   per-sample training behaviour; \code{DropPath(0)} byte-for-byte
   equality with the no-op; attention pool shape; state-dict key
   invariance under \cfgkey{drop\_path\_rate}; \code{attn\_pool.*} key
   appearance.

\item[\file{tests/pipelines/test\_submit\_tta.py} (5 tests)]
   Phase~2: auto-derived left/right class permutation; no-TTA path;
   untrained-mask path; flip-TTA softmax reinforcement of the correct
   class after remap.

\item[\file{tests/shared/data/test\_still\_frames\_dataset.py} (6 tests)]
   Phase~3 (DINO, retained): frame-path discovery across both folder
   layouts; deduplication across roots; multi-view sampling shape.

\item[\file{tests/shared/models/test\_dino\_ssl.py} (4 tests)]
   Phase~3 (DINO, retained): SSL trunk key superset of the supervised
   backbone keys; finite back-propagating loss; teacher-EMA midpoint
   identity.

\item[\file{tests/pipelines/test\_init\_from\_ssl.py} (1 test)]
   Phase~3 (DINO, retained) end-to-end smoke test for the
   \cfgkey{model.init\_from} hook.

\item[\file{tests/shared/models/test\_vjepa\_ssl.py} (10 tests)]
   Phase~3 (V-JEPA): trunk forward shape; remapped V-JEPA trunk keys
   cover the supervised backbone \emph{exactly} (zero missing, zero
   unexpected); \code{load\_state\_dict} only misses the classifier;
   the L1 loss is zero for an empty mask (with a valid \code{grad\_fn})
   and finite-and-backprop-able otherwise; \code{apply\_frame\_mask}
   zeroes the correct positions; mask sampler bounds and
   ``never-all-masked'' guarantee; teacher EMA midpoint.

\item[\file{tests/shared/data/test\_clip\_ssl\_dataset.py} (7 tests)]
   Phase~3 (V-JEPA): both folder layouts; cross-root deduplication;
   missing-root skip; sample shape; temporal jitter under hold-out.

\item[\file{tests/pipelines/test\_init\_from\_vjepa.py} (1 test)]
   Phase~3 (V-JEPA) end-to-end smoke test that explicitly pins the
   no-double-\code{backbone.}-prefix invariant of the V-JEPA save format.

\item[\file{tests/shared/utils/test\_class\_balance.py},\\
      \file{tests/shared/data/test\_time\_reversal.py},\\
      \file{tests/pipelines/test\_submit\_phase4\_tta.py}]
   Phase~4: per-class weight policies (\code{none}, \code{inverse},
   \code{sqrt\_inverse}, \code{cb}) and their corner cases (zero-count
   class set to weight 0); time-reversal table-building and
   label-remap symmetry (pair-permuted twice = identity); multi-scale
   TTA shape and logit-adjustment numerics.
\end{description}

% =====================================================================
\section{Expected Impact}
\label{sec:strategy}

\subsection{Per-change order-of-magnitude estimates}

Based on the literature and on the gap currently observed between
internal validation and the public leaderboard, we expect the following
contributions to Track~A's public Top-1 score (these are
order-of-magnitude estimates, not promises):

\begin{center}\small
\begin{tabular}{lcc}
\toprule
\textbf{Change} & \textbf{Internal val} & \textbf{Public Top-1} \\
\midrule
Baseline (current best)                        & 0.5666 & 0.3584 \\
+ official val split (Phase~1)                 & lower\,$^*$ & --     \\
+ Kaiming init (Phase~1)                       & +0.5--2 pt & +0.3--1 pt \\
+ AMP (Phase~1)                                & speedup only & speedup only \\
+ optimizer/scheduler state on resume (Phase~1)& 0--1 pt & 0--1 pt \\
+ untrained-class masking (Phase~1)            & negligible & +0.5--1.5 pt \\
+ DropPath (Phase~2)                           & +0.5--1 pt & +0.3--0.8 pt \\
+ Attention head (Phase~2)                     & +0.5--1 pt & +0.3--0.8 pt \\
+ EMA model averaging (Phase~2)                & +0.3--0.8 pt & +0.2--0.6 pt \\
+ TTA flip with class-pair remap (Phase~2)     & +0--0.5 pt & +0.5--1 pt \\
+ AdamW + WD~0.05 (Phase~2)                    & +0.3--1 pt & +0.3--1 pt \\
+ DINO SSL pretraining (deprecated)            & --1 to +1 pt\,$^\dagger$ & --1 to +1 pt\,$^\dagger$ \\
+ V-JEPA SSL pretraining (Phase~3)             & +2--4 pt & +2--4 pt \\
+ Class-balanced sampler + CB loss (Phase~4)   & +1--2 pt & +1--2 pt \\
+ Time-reversal augmentation (Phase~4)         & +0.5--1.5 pt & +0.5--1.5 pt \\
+ Multi-scale TTA (Phase~4)                    & +0.3--0.8 pt & +0.3--0.8 pt \\
+ Logit adjustment (Phase~4)                   & +0.2--0.5 pt & +0.2--0.5 pt \\
\bottomrule
\end{tabular}\\
\smallskip
{\footnotesize $^*$ The official split is harder than the leaky
internal split, so reported val will drop --- but the public score will
move closer to the val score.\\
$^\dagger$ As observed in practice and explained in
Section~\ref{sec:phase3:dino-deprecated}: DINO on still frames produces
inconsistent downstream gains. We retain the row in this table only as
a baseline against which V-JEPA's stable gains stand out.}
\end{center}

% =====================================================================
\section{Legacy DINO Pretraining (Deprecated)}
\label{sec:phase3dino}

The single highest-impact change available within Track~A is to give the
trunk something useful to do with the unlabeled frames before the
supervised loss starts back-propagating. We implemented DINO-v1
\cite{Caron2021DINO}, a teacher-student knowledge-distillation framework
that has consistently outperformed contrastive methods (MoCo, SimCLR) on
ResNet-50 trunks at the dataset sizes we are working with.

\subsection{Track-A compliance}

The competition forbids \emph{external data} and \emph{pretrained
models}. Self-supervised pretraining on the \emph{provided} unlabeled
images is therefore allowed: the labels are simply not consumed and the
trunk starts from random initialisation. We pretrain on the union of the
\file{train/}, \file{val/} and \file{test/} folders, which gives roughly
\(\sim 12{,}000\) still frames (\(\sim 3{,}000\) videos
\(\times\) 4 frames) and is sufficient to drive the DINO loss for many
hundreds of steps per epoch at batch size~64.

\subsection{Still-frames dataset}
\label{sec:phase3dino:dataset}

A new file \file{src/smth2smth/shared/data/still\_frames\_dataset.py}
exposes two utilities:

\begin{itemize}[leftmargin=2em]
\item \code{collect\_all\_frame\_paths(roots)} walks the provided
      directories. It transparently handles both layouts in use:
      \file{train/<class>/<video>/<frame>.jpg} (via
      \code{collect\_video\_samples}) and
      \file{test/<video>/<frame>.jpg} (one level shallower, no class
      bucket). Duplicate paths are removed.
\item \code{MultiViewStillFramesDataset} returns, for each frame, two
      global views (typically 224\(\times\)224 with random crop scale in
      \([0.4, 1.0]\)) and \(L\) local views (typically 96\(\times\)96
      with random crop scale in \([0.05, 0.4]\)). The augmentations
      include horizontal flip, colour jitter, grayscale, and Gaussian
      blur, following the DINO recipe.
\end{itemize}

The dataset returns a \code{dict} with keys \code{global\_views} and
\code{local\_views}; a custom collate function stacks them into per-view
batches.

\subsection{The DINO model}
\label{sec:phase3dino:model}

\file{src/smth2smth/shared/models/dino\_ssl.py} contains three pieces.

\paragraph{Trunk.} \code{build\_resnet50\_trunk()} returns a
\code{torchvision.models.resnet50(weights=None)} with its final \code{fc}
replaced by \code{nn.Identity}. The trunk emits 2048-d features.

\paragraph{Head.} \code{DinoHead} is the canonical 3-layer projection
MLP used in DINO: \code{Linear} $\to$ \code{GELU} $\to$ \code{Linear}
$\to$ \code{GELU} $\to$ \code{Linear} (to a low-dimensional bottleneck)
$\to$ L2 normalisation $\to$ weight-normalised \code{Linear} mapping to
\(K\) prototypes (\(K=4096\) by default).

\paragraph{Composite.} \code{DinoModel} composes the trunk and the head.
Two copies are instantiated at training time --- a \emph{student} and a
\emph{teacher}. Both have identical architecture; the teacher's weights
are an EMA of the student's and are never updated by gradient descent.

\subsection{The DINO loss}
\label{sec:phase3dino:loss}

Let \(\mathcal{V} = \{v_1, \ldots, v_{2+L}\}\) be the set of all views
of one image (two globals, \(L\) locals). Let \(p_s(v)\) and \(p_t(v)\)
be the softmax-normalised outputs of the student and teacher at
temperatures \(\tau_s\) and \(\tau_t\) respectively. The loss is
\[
\mathcal{L} \;=\;
\frac{1}{|\mathcal{P}|}\!\!\sum_{(t, s)\in \mathcal{P}}\!\!
H\!\left( p_t(v_t),\; p_s(v_s) \right),
\]
where \(\mathcal{P}\) is the set of (teacher view, student view) pairs
with \emph{distinct} indices and \(H\) is the standard cross-entropy.
The teacher only processes the two global views (so
\(|\mathcal{P}| = 2 \cdot |\mathcal{V}| - 2 \cdot 1\)), but the student
processes \emph{all} views. The teacher's logits are sharpened
(\(\tau_t = 0.04\) is much smaller than \(\tau_s = 0.1\)) and centred by
an EMA of the batch mean to prevent collapse:
\[
c^{(t+1)} \;=\; m_c\, c^{(t)} \;+\;
(1 - m_c)\, \frac{1}{B}\sum_{i=1}^{B}\mathrm{logits}_t(v_i^{\mathrm{global}}),
\]
with momentum \(m_c = 0.9\).

\begin{lstlisting}[style=py]
class DinoLoss(nn.Module):
    def forward(self, student_logits_per_view, teacher_logits_per_global):
        teacher_probs_per_global = [
            F.softmax((logits - self.center) / self.teacher_temp,
                      dim=-1).detach()
            for logits in teacher_logits_per_global]
        student_log_softmax_per_view = [
            F.log_softmax(logits / self.student_temp, dim=-1)
            for logits in student_logits_per_view]
        total_loss, n_terms = 0.0, 0
        for t_idx, t_probs in enumerate(teacher_probs_per_global):
            for s_idx, s_log in enumerate(student_log_softmax_per_view):
                if t_idx == s_idx:
                    continue
                total_loss = total_loss \
                             + (-(t_probs * s_log).sum(dim=-1).mean())
                n_terms += 1
        # ... centering buffer EMA update ...
        return total_loss / n_terms
\end{lstlisting}

\subsection{Teacher EMA momentum schedule}

The teacher's parameters are updated by
\(\theta_t \leftarrow m\,\theta_t + (1-m)\,\theta_s\) after every
optimizer step, with \(m\) following a cosine schedule from
\(m_0 = 0.996\) at the start to \(m_1 = 1.0\) at the end of training.
This mirrors the recipe in the DINO paper and stabilises training
significantly --- a constant momentum is more prone to collapse on small
datasets.

\subsection{Trunk-only checkpoint format}

After every epoch we save a small ($\approx$ 95~MiB) checkpoint that
contains \emph{only} the trunk parameters, with the \code{trunk.}
prefix stripped:

\begin{lstlisting}[style=py]
trunk_state_dict = {
    k.removeprefix("trunk."): v
    for k, v in student.state_dict().items()
    if k.startswith("trunk.")
}
torch.save({"trunk_state_dict": trunk_state_dict,
            "epoch": epoch + 1}, out_path)
\end{lstlisting}

This format is intentionally minimal: there is no point keeping the DINO
head, the teacher copy or the optimizer state once supervised training
takes over.

\subsection{The \cfgkey{model.init\_from} hook in \file{train.py}}
\label{sec:phase3dino:initfrom}

The supervised trainer received a single new opt-in:

\begin{lstlisting}[style=py]
init_from = cfg.model.get("init_from") if hasattr(cfg.model, "get") \
            else None
if init_from:
    payload = torch.load(Path(str(init_from)).resolve(),
                         map_location=device, weights_only=False)
    trunk_state = payload.get("trunk_state_dict")
    prefixed = _ssl_trunk_to_supervised_keys(trunk_state)
    missing, unexpected = model.load_state_dict(prefixed, strict=False)
\end{lstlisting}

The non-trivial part is the key remap. The supervised
\code{AvancedResNet50TSM} wraps every bottleneck block's \code{conv1}
in
\[
\code{block.conv1} \;\leftarrow\;
\code{nn.Sequential(TemporalShift(\(T\), shift\_div), conv1)},
\]
which means the parameter that the SSL trunk has at key
\code{layer1.0.conv1.weight} now lives at
\code{layer1.0.conv1.1.weight} in the supervised model.
\code{\_ssl\_trunk\_to\_supervised\_keys} performs this rename
unconditionally for every key matching
\code{layer[1-4]\textbackslash.\textbackslash d+\textbackslash.conv1\textbackslash.(weight|bias)}:

\begin{lstlisting}[style=py]
block_conv1_re = re.compile(
    r"^(layer[1-4]\.\d+\.conv1)\.(weight|bias)$")
out: dict[str, torch.Tensor] = {}
for k, v in trunk_state.items():
    match = block_conv1_re.match(k)
    if match is not None:
        new_k = f"{match.group(1)}.1.{match.group(2)}"
    else:
        new_k = k
    out[f"backbone.{new_k}"] = v
return out
\end{lstlisting}

The load is intentionally non-strict so that the supervised model's
classifier head and (when enabled) attention pool keys remain at their
fresh-init values. The trainer logs how many tensors were loaded, and
warns loudly if any backbone key was \emph{not} covered by the SSL
checkpoint.

\subsection{Dual-stream RGB + frame-difference TSM}
\label{sec:dual-stream-rgb-diff}

Beyond the single-stream \code{AvancedResNet50TSM} baseline, the codebase
includes \code{dual\_stream\_rgb\_diff\_tsm}: two backbones with TSM
inserted in the usual bottleneck (ResNet-50 on RGB frames) and
basic-block (ResNet-34 on consecutive RGB differences
\(I_t - I_{t-1}\)) branches. Per time index \(t=0,\ldots,T-1\) the model
forms a fused token by concatenating the RGB embedding at \(t\) with a
motion embedding---zero-padded at \(t{=}0\), otherwise the difference
embedding for \((t{-}1,t)\)---and projecting to a fixed width
(\code{model.fuse\_dim}, default~512). A shared temporal readout (mean or
1-query attention, matching the single-stream head options) and linear
classifier produce clip logits. See
\file{docs/diagrams/dual\_stream\_rgb\_diff\_attention.svg}
for a block diagram of the architecture.

\paragraph{Configuration.} Model defaults live in
\file{configs/model/dual\_stream\_rgb\_diff\_tsm.yaml}. A ready-made
Track~A training recipe (30~epochs, warmup, RandAugment-T + FrameMixUp,
official validation split, AMP, EMA, class boosting with deterministic
paired-verb duplicate rows, validation every epoch via
\code{training.eval\_every\_n\_epochs=1}) is
\file{configs/experiment/track\_a\_dual\_stream\_30e\_class\_boost.yaml}.
Stochastic \code{temporal\_reversal\_augment} is disabled automatically
when \code{dataset.class\_boosting.enabled} is true. The README documents
the Hydra invocation.

\paragraph{SSL warm-start.} The \code{model.init\_from} hook from
\S\ref{sec:phase3:initfrom} applies only to \code{AvancedResNet50TSM}
backbone key names; training with \code{dual\_stream\_rgb\_diff\_tsm}
skips loading so that mis-matched prefixes are not applied silently.

\subsection{Per-backbone learning rates for the dual-stream model}
\label{sec:dual-stream-per-backbone-lr}

Empirically, training the RGB ResNet-50 and the motion ResNet-34 with the
\emph{same} learning rate is suboptimal: the motion branch sees a noisier
input signal (consecutive-frame difference \(I_t-I_{t-1}\) rather than
raw RGB), has fewer parameters, and tends to overfit before the deeper
RGB trunk has converged. Two new keys in \file{configs/train/default.yaml}
expose a per-branch peak LR without changing legacy behaviour:

\begin{itemize}[leftmargin=2em]
\item \cfgkey{training.motion\_lr} --- absolute peak LR for the
   ResNet-34 motion branch.  Overrides the ratio when both are set.
\item \cfgkey{training.motion\_lr\_ratio} --- relative multiplier; the
   motion branch peak LR is \(\text{lr}\times\text{ratio}\).  Default
   \code{null} keeps the historical single-LR-group path.
\end{itemize}

When either knob is set and the model is the dual-stream builder,
\code{pipelines/train.py} constructs three optimiser groups
(\code{rgb\_backbone}, \code{motion\_backbone}, \code{fuse\_head}) with
per-group \code{warmup\_lr\_max}, so linear warmup and cosine annealing
both respect the per-branch peak LR.  The shared \cfgkey{min\_lr} is
reused as the cosine \(\eta_{\min}\) for all groups.  This path is
mutually exclusive with the LoRA two-group path (LoRA is not used by
the dual-stream model, so the choice is well-defined).

\paragraph{Long-run preset.}  The companion experiment YAML
\file{configs/experiment/track\_a\_dual\_stream\_long\_class\_boost.yaml}
extends the 30-epoch class-boost preset with: an 80-epoch budget so
cosine has room to anneal, \code{motion\_lr\_ratio: 0.3}, aggressive
early stopping (\code{early\_stopping\_patience: 3}, \(\Delta=0.001\)),
and a separate checkpoint path.  The bash launcher
\file{scripts/run\_dual\_stream\_long\_with\_fallback.sh} chains a
slower-LR fallback: if Run~1's log contains
\code{Early stopping triggered}, it launches a second run with
\code{lr=0.6\,\(\times\)\,lr$_1$} (default \code{6e-5}) into a fresh
log/checkpoint pair, otherwise it exits cleanly.

\paragraph{Result.}  Run~A1 was stopped externally after epoch~16
(31.79\,\% EMA val top-1).  Run~A1b resumed from the \texttt{*.last.pt}
checkpoint with a 20\,\% higher RGB/fuse LR (\texttt{lr}=1.2e-4,
\texttt{resume\_apply\_cfg\_lr}=1), trained through epoch~39, and
triggered early stopping (patience~3).  Best EMA validation top-1:
\(\mathbf{39.47\,\%}\) (\texttt{ckpt}: \file{.../dual\_stream\_long\_class\_boost\_lr1\_20260515\_025557.pt}).
Flip-TTA CSV \file{.../track\_a\_dual\_stream\_long\_lr1\_resume\_20260515.csv}
(\(6\,913\) test clips) scored \(\mathbf{37.69\,\%}\) public top-1 on
the Kaggle leaderboard---\(\sim 1.9\)\,pp above the previous Track-A
best (35.84\,\%, Tab.~in \S\ref{sec:context}) and within \(\sim 1.8\)\,pp
of the official validation proxy (39.47\,\%), consistent with the
official-val fix from Phase~1.

\subsection{Track A symbol glossary (training runs)}
\label{sec:tracka:glossary}

Mirrors the Track-B convention from \S\ref{sec:trackb:glossary}; symbols
are reused where they exist in both tracks so a reader does not have to
juggle two conventions.  Compact assignments below are the ones used in
the Track-A run logbook (Tab.~\ref{tab:tracka-runlog}).

\begin{table}[ht]
\centering\footnotesize
\begin{tabular}{@{}r p{0.36\linewidth} p{0.48\linewidth}@{}}
\toprule
\textbf{Sym.} & \textbf{Hydra key(s)} & \textbf{Meaning} \\
\midrule
\texttt{expt} & \texttt{experiment=\dots} & Preset YAML under
\file{configs/experiment/}. \\
\texttt{T} & \texttt{dataset.num\_frames} & Frames per clip. \\
\texttt{Sz} & \texttt{dataset.image\_size} & Square crop side. \\
\texttt{bs} & \texttt{training.batch\_size} & Clips per optimiser step. \\
\texttt{E} & \texttt{training.epochs} & Total epoch budget. \\
\texttt{lr} & \texttt{training.lr} & Base / RGB-branch peak LR (AdamW). \\
\texttt{m\_r} & \texttt{training.motion\_lr\_ratio} & Motion-branch
peak LR as a fraction of \(\text{lr}\); \code{null} for legacy single-group. \\
\texttt{wd} & \texttt{training.weight\_decay} & AdamW weight decay. \\
\texttt{wu} & \texttt{training.warmup\_epochs} & Linear LR warm-up length. \\
\texttt{cos} & \texttt{training.scheduler\_cosine} & Cosine anneal after warm-up. \\
\texttt{ls} & \texttt{training.label\_smoothing} & Soft-label smoothing for CE. \\
\texttt{amp} & \texttt{training.amp} & Mixed precision. \\
\texttt{ema} & \texttt{training.ema\_enabled} & EMA of weights; decay in \texttt{ema\_decay}. \\
\texttt{es} & \texttt{training.early\_stopping\_patience} & Patience (val top-1 evals without improvement) before early stop. \\
\texttt{cb} & \texttt{dataset.class\_boosting.enabled} & Deterministic paired-verb row duplication. \\
\texttt{aug} & augment group & Spatial/temporal augment preset. \\
\texttt{ckpt} & \texttt{training.checkpoint\_path} & Primary weights file. \\
\bottomrule
\end{tabular}
\caption{Symbols used in Tab.~\ref{tab:tracka-runlog}.  Extra CLI
overrides are appended verbatim after a semicolon.}
\label{tab:tracka-glossary}
\end{table}

\subsection{Track A training run logbook}
\label{sec:tracka:runlog}

Statuses: \textbf{RUNNING} (PID alive or log still progressing),
\textbf{DONE} (trainer printed \texttt{Done.}), \textbf{FAILED}
(traceback or non-zero exit).  Update the corresponding row when a job
finishes; add new rows for every fresh \texttt{nohup} launch.

\begin{table}[ht]
\caption{Chronological log of substantial Track-A training jobs.}
\label{tab:tracka-runlog}
\begin{minipage}{\linewidth}\small\sloppy
\begin{description}[leftmargin=1.5em, style=nextline]
\item[\textbf{Run A1} (STOPPED)]
  Wrapper: \file{scripts/run\_dual\_stream\_long\_with\_fallback.sh}.
  Log: \texttt{logs/track\_a\_dual\_stream\_long\_lr1\_20260515\_025557.log}.
  Preset \texttt{track\_a\_dual\_stream\_long\_class\_boost}
  (Tab.~\ref{tab:tracka-glossary}); \texttt{aug}=randaugment\_t;
  \texttt{T}=4; \texttt{Sz}=224; \texttt{bs}=16;
  \texttt{E}=80; \texttt{lr}=1e-4; \texttt{m\_r}=0.3
  (\(\Rightarrow\) motion lr=3e-5);
  \texttt{wd}=0.05; \texttt{wu}=6; \texttt{cos}=1; \texttt{ls}=0.1;
  \texttt{amp}=1; \texttt{ema}=0.999;
  \texttt{es}=3 (\(\Delta=0.001\)); \texttt{cb}=1.
  Stopped externally mid epoch~17.  Best EMA val top-1 (epoch~16):
  \textbf{31.79\,\%}.
\item[\textbf{Run A1b} (DONE)]
  Log: \texttt{logs/track\_a\_dual\_stream\_long\_lr1\_resume\_20260515\_104715.log}.
  Resume from \texttt{*.last.pt} (epoch~16); override
  \texttt{lr}=1.2e-4 (+20\%), \texttt{m\_r}=0.3,
  \texttt{resume\_apply\_cfg\_lr}=1; same preset otherwise.
  Early stop at epoch~39 (patience exhausted).  Best EMA val top-1:
  \textbf{39.47\,\%}; same \texttt{ckpt} as Run~A1.
  Submit (flip TTA): \texttt{.../track\_a\_dual\_stream\_long\_lr1\_resume\_20260515.csv};
  log \texttt{.../track\_a\_dual\_stream\_long\_lr1\_resume\_20260515\_submit.log}.
  Public Kaggle top-1: \textbf{37.69\,\%}.
\item[\textbf{Run A2 / truite} (DONE -- FAILED LEARNING)]
  Log: \texttt{logs/truite.log}.
  Preset \texttt{track\_a\_vit\_truite}: video\_mae\_vit \texttt{variant}=vit\_b,
  \texttt{head}=\textbf{mean} (mean pool, MP), \(\sim\)87\,M params;
  same train recipe as the fish-matrix presets below.
  \texttt{T}=4; \texttt{Sz}=224; \texttt{bs}=16; \texttt{E}=40;
  \texttt{lr}=5e-4; \texttt{wu}=5; RandAug-T (\texttt{n}=4, \texttt{m}=7);
  cube CutMix; \texttt{cb}=0.  Early stop after 18 epochs; best val top-1
  \textbf{5.41\,\%}; train top-1 plateaued near \(\sim 7\)\,\%
  (loss \(\ln 33 \approx 3.5\)).
\item[\textbf{Run A3 / sardine} (STOPPED -- FAILED LEARNING)]
  Log: \texttt{logs/track\_a\_vit\_sardine\_20260515\_230200.log}.
  Preset \texttt{track\_a\_vit\_sardine}; \texttt{variant}=vit\_s; MP;
  \(\sim\)22\,M; \texttt{bs}=64; \texttt{lr}=1e-3; \texttt{wu}=10;
  RandAug-T (\texttt{n}=6, \texttt{m}=10).  Stopped manually after epoch~3:
  \(\sim 6.7\,\%\) train / \textbf{6.15\,\%} best val (epoch~2)---same
  failure mode as truite despite 4\(\times\) larger batch and heavier aug.
\item[\textbf{Run A4 / truite AP} (STOPPED -- FAILED LEARNING)]
  Log: \texttt{logs/track\_a\_vit\_truite\_ap\_20260515\_234500.log}.
  Preset \texttt{track\_a\_vit\_truite\_ap}: ViT-B with \textbf{attentive probe}
  (\texttt{head=attn}, 4 heads), \(\sim\)89\,M params; otherwise identical to
  truite.  Stopped manually mid-training (same plateau as MP): e.g.\ mid-epoch
  steps still show train loss \(\approx 3.35\) and train top-1 \(\approx
  6.5\)--\(7\)\,\%---\emph{no gain} from replacing mean pool.  Confirms the
  bottleneck is not the temporal readout.
\item[\textbf{Run A5 SSL / truite (slot 9)} (PRETRAIN DONE; FT NOT RUN)]
  Phase~1 VideoMAE SSL (\texttt{smth2smth.pipelines.pretrain\_videomae});
  preset \texttt{track\_a\_ssl\_pretrain\_truite} (ViT-S baseline, no
  \file{SSL.md} \S7 overrides).  Log: \texttt{logs/ssl\_truite\_pretrain.log}.
  Launcher: \texttt{.venv/bin/python} with \texttt{PYTHONPATH=src}
  (\texttt{nohup}; minimal shells often lack \texttt{uv} on \texttt{PATH}).
  MAE \texttt{E}=100, \texttt{wu}=5, \texttt{bs}=16, \texttt{lr}=1.5e-4,
  \texttt{mask}=0.75; \(\sim 58\,651\) unlabeled clips; avg loss
  \(0.851 \to 0.256\) (ep~100); encoder
  \file{checkpoints/track\_a/ssl/truite\_encoder.pt} (last epoch, rewritten
  each epoch).  \textbf{Phase~2 not run:} no \texttt{ssl\_truite\_finetune.log},
  no \texttt{truite\_ft.pt} (pending \texttt{track\_a\_ssl\_finetune\_truite}
  after encoder).  Distinct from supervised Runs~A2/A4 (\texttt{track\_a\_vit\_*}).
\item[\textbf{Run rouget} (ViT-B + MP + class boosting; DONE -- FAILED)]
  Preset \texttt{track\_a\_vit\_rouget}: truite stack + deterministic class
  boosting (\texttt{cb}=1).  Representative log (epoch~16): train
  \(\sim 8.6\,\%\) top-1, val \(\sim 4.6\,\%\), best val still only
  \textbf{5.59\,\%} with patience counting down---no useful convergence.
\item[\textbf{Run roussette} (ViT-B + MP + class balancing; DONE -- FAILED)]
  Preset \texttt{track\_a\_vit\_roussette}: sqrt-inverse sampler + CB loss,
  no boosting.  Early stop at epoch~23; best val top-1 \textbf{5.23\,\%}.
\item[\textbf{Run raie} (ViT-B + MP + boosting + balancing; DONE -- FAILED)]
  Preset \texttt{track\_a\_vit\_raie}: both \texttt{cb} and balancing.
  Best val top-1 \textbf{6.00\,\%}; example epoch~15: val \(\sim
  6.1\,\%\) live / \(\sim 5.9\,\%\) EMA with train top-1 still
  \(\sim 5.7\)--\(5.8\,\%\)---long-tail tricks do not substitute for
  convolutional inductive bias or SSL features.
\item[\textbf{SSL slot 8 / barbeau} (NOT RUN)]
  Thomas slot: ViT-B, MAE \texttt{warmup\_epochs}=10
  (\texttt{track\_a\_ssl\_pretrain\_barbeau} /
  \texttt{track\_a\_ssl\_finetune\_barbeau}).  Job was not launched.
\item[\textbf{SSL slot 12 / raie} (DONE)]
  Romain slot: attentive probe \texttt{head\_num\_heads}=6.
  \paragraph{Phase 1.} 100/100 MAE after one AMP \texttt{scatter\_} crash
  (retry); loss \(0.851 \to 0.254\);
  \file{raie\_encoder.pt}; logs \texttt{ssl\_raie\_pretrain.log} /
  \texttt{.crash.log}.
  \paragraph{Phase 2.} 30/30 via \texttt{chain\_raie.sh}; official val only;
  train top-1 \textbf{44.15\,\%}; best EMA val \textbf{34.80\,\%} (ep~29,
  \file{raie\_ft.pt}); log \texttt{ssl\_raie\_finetune.log}.
  Failed post-sweep resume (OOM) excluded.  No submission.
\item[\textbf{SSL slot 10 / roussette} (DONE)]
  Romain slot: YAML \texttt{seed}=123 (§7); resolved run used \texttt{seed}=42;
  pretrain \texttt{bs}=8 vs planned 16 (CLI).
  \paragraph{Phase 1.} 100/100 MAE; loss \(0.843 \to 0.251\);
  \file{checkpoints/track\_a/ssl/roussette\_encoder.pt};
  log \texttt{logs/ssl\_roussette\_pretrain.log}.
  \paragraph{Phase 2.} 30/30 FT; train top-1 \textbf{41.44\,\%} (ep~30);
  best EMA val top-1 \textbf{33.61\,\%} (ep~27, saved);
  final val \textbf{33.62\,\%} live / \textbf{33.31\,\%} EMA;
  \file{roussette\_ft.pt}; log \texttt{logs/ssl\_roussette\_finetune.log}.
  No submission yet.
\item[\textbf{SSL slot 16 / murene} (NOT RUN --- dataset quota)]
  Romain slot: FT \texttt{lr}=3e-4
  (\texttt{track\_a\_ssl\_pretrain\_murene} /
  \texttt{track\_a\_ssl\_finetune\_murene}).  Job never started: dataset
  download failed (\textbf{user limit exceeded in 24h} on the VM).  No
  \texttt{ssl\_murene\_*} logs or checkpoints.
\item[\textbf{SSL slot 15 / piranha} (DONE)]
  Romain slot: FT \texttt{epochs}=40 (vs 30 default).
  \paragraph{Phase 1.} 100/100 MAE; final loss \textbf{0.2528};
  \file{piranha\_encoder.pt}; log \texttt{logs/ssl\_piranha\_pretrain.log}.
  \paragraph{Phase 2.} 40/40; train top-1 \textbf{49.54\,\%} (ep~40);
  best EMA val \textbf{35.60\,\%} (ep~32, \file{piranha\_ft.pt});
  log \texttt{logs/ssl\_piranha\_finetune.log}.  Chained launcher vs
  two-step \file{SSL.md}; Hydra otherwise matches preset.  No submission.
\item[\textbf{SSL slot 14 / thon} (NOT RUN --- this host)]
  Romain slot: MAE \texttt{mask\_ratio}=0.85
  (\texttt{track\_a\_ssl\_pretrain\_thon} /
  \texttt{track\_a\_ssl\_finetune\_thon}).  No
  \texttt{logs/ssl\_thon\_*} or \file{thon\_encoder.pt} /
  \file{thon\_ft.pt}; presets match \file{SSL.md} but job never
  launched on this VM (no partial log or \texttt{.pid}).  Config drift N/A.
  Superseded by overnight batch \textbf{E5} below (same VM, new presets).
\item[\textbf{E5 / thon batch} (OVERNIGHT 20260516)]
  VM \textbf{thon}; honest val; ViT TTA scales
  \texttt{[0.857, 1.0, 1.143]} at submit.
  \paragraph{Phase 1 MAE (DONE).}
  Preset \texttt{track\_a\_ssl\_pretrain\_e5};
  log \texttt{logs/thon\_e5\_mae\_pretrain\_20260516.log};
  pid \texttt{logs/thon\_e5\_mae\_pretrain\_20260516.pid}.
  \texttt{E}=100; \texttt{wu}=5; \texttt{bs}=16; \texttt{lr}=1.5e-4;
  \texttt{mask}=0.75; minimal aug.  MAE loss \(0.856 \to 0.258\) (ep~100);
  encoder \file{checkpoints/track\_a/ssl/e5\_encoder.pt} (85\,MiB).
  \paragraph{Phase 2a FT repr.\ (DONE).}
  Preset \texttt{track\_a\_ssl\_finetune\_e5}; \texttt{bs}=8; \texttt{E}=50;
  \texttt{head}=attn (4 heads).  Log
  \texttt{logs/thon\_e5\_ft\_stage1\_rep50\_20260516.log} (VM restart ep~24,
  resumed ep~25--50).  Best EMA val \textbf{37.60\,\%} (ep~45);
  final ep~50 EMA \textbf{37.32\,\%}; \file{e5\_ft.pt}.
  \paragraph{Phase 2a cont.\ (RUNNING).}
  Preset \texttt{track\_a\_ssl\_finetune\_e5\_continue150\_bs16};
  \texttt{E}=150; \texttt{bs}=16; cont.\ rewarm ratio=0.5, \texttt{wu}=6.
  Log \texttt{logs/thon\_e5\_ft\_continue150\_bs16\_20260516.log};
  pid \texttt{logs/thon\_e5\_ft\_continue150\_bs16\_20260516.pid};
  resume \file{e5\_ft.last.pt} (ep~50).  Phase~2b cRT pending.
  \item[\textbf{SSL slot 13 / sole} (NOT RUN --- this host)]
  Romain slot: attentive probe \texttt{head\_num\_heads}=8
  (\texttt{track\_a\_ssl\_pretrain\_sole} /
  \texttt{track\_a\_ssl\_finetune\_sole}).  No
  \texttt{logs/ssl\_sole\_pretrain.log} or \texttt{ssl\_sole\_finetune.log};
  no \file{sole\_encoder.pt} / \file{sole\_ft.pt} under
  \file{checkpoints/track\_a/ssl/} on this machine (launch deferred while
  GPU was occupied).  Presets on disk match \file{SSL.md}; no config drift
  logged.  Distinct from supervised \texttt{track\_a\_vit\_*} mean-pool jobs.
\item[\textbf{SSL slot 11 / rouget} (DONE)]
  VideoMAE SSL sweep (see \file{SSL.md}): ViT-S, long MAE
  (\texttt{pretrain.epochs}=200), attentive probe (4 heads).
  \paragraph{Phase 1.} Pretrain on train+val+test unlabeled clips;
  encoder \file{checkpoints/track\_a/ssl/rouget\_encoder.pt};
  log \texttt{logs/ssl\_rouget\_pretrain.log}.
  \paragraph{Phase 2a (honest val).} \texttt{track\_a\_ssl\_finetune\_rouget},
  train=\texttt{data/train} only, official \texttt{data/val} for metrics;
  30 epochs; best EMA val top-1 \textbf{34.96\,\%};
  log \texttt{logs/ssl\_rouget\_finetune.log}.
  \paragraph{Phase 2b (invalid val).} Resumed with
  \texttt{dataset.include\_val\_in\_train=true}; stopped at epoch~44/60;
  val top-1 rose to \(\sim 60\,\%\) (\emph{leaky}---val labels in training).
  Log \texttt{logs/ssl\_rouget\_finetune\_trainval.log}.
  \paragraph{Submit.} \file{submissions/track\_a\_rouget\_ssl.csv} from
  post--Phase-2b \texttt{rouget\_ft.pt}; flip-only TTA
  (\texttt{tta\_scales=[1.0]}).  Public Kaggle top-1: \textbf{36.0\,\%}---
  aligns with Phase~2a val, not the inflated \(\sim 60\,\%\); training on
  val without a held-out set likely \textbf{overfit} and did not beat the
  dual-stream best (\textbf{37.69\,\%}, Run~A1b).
\item[\textbf{Run E4 / truite batch 20260516} (RUNNING)]
  Overnight batch \file{instruction\_16052026.md}: supervised R50+TSM from
  scratch (no SSL).  Preset \texttt{track\_a\_e4\_tsm\_sgdr};
  \texttt{expt}=E4; \texttt{T}=4; \texttt{Sz}=224; \texttt{bs}=16;
  \texttt{E}=90; \texttt{lr}=1e-4; \texttt{wd}=0.05;
  scheduler SGDR \texttt{T\_0}=30 (snapshots at ep~30/60/90);
  \texttt{aug}=randaugment\_t; \texttt{ls}=0.1; \texttt{amp}=1;
  \texttt{ema}=0.999; frame mixup; \texttt{tta} scales
  \([0.875,1.0,1.125]\) (CNN).  Honest official val only.
  Log \texttt{logs/truite\_e4\_supervised\_sgdr90\_20260516.log};
  PID \texttt{logs/truite\_e4\_supervised\_sgdr90\_20260516\_resume.pid}
  (VM restart 2026-05-17; resumed from \texttt{e4\_ft.last.pt} at ep~32).
  \texttt{ckpt}=\file{checkpoints/track\_a/e4\_ft.pt} (+ \texttt{e4\_ft\_snap1.pt};
  cycles 2--3 pending).  Interim best val top-1 \textbf{38.67\,\%} (EMA, ep~50);
  \texttt{e4\_ft\_snap1.pt} at cycle~1.  \textbf{Phase~2 prepared:}
  \texttt{track\_a\_e4\_tsm\_continue150} (ep~91--150, \texttt{wu}=10, cosine
  \texttt{T\_max}=50, \texttt{bs}=16); launch after 90-ep \texttt{Done.}
  Log \texttt{logs/truite\_e4\_continue150\_20260516.log};
  chain watcher \texttt{logs/truite\_e4\_chain\_continue150\_20260516.log}
  (auto-starts phase~2 when 90-ep job exits).
  \textbf{Phase~3 trainval FT} (STOPPED): \texttt{track\_a\_e4\_tsm\_trainval\_ft};
  train+val; submit at ep~6 \texttt{track\_a\_e4\_trainval\_ep6\_20260516.csv}:
  Kaggle public top-1 \textbf{42.92\,\%}.  Training halted at ep~17/25 (leaky val
  $\gg$ LB).  Log \texttt{logs/truite\_e4\_trainval\_ft\_20260516.log}.
  Honest-val best remains \texttt{e4\_ft.pt} (43.31\,\% EMA, ep~109).
\item[\textbf{Batch E6 / roussette} (RUNNING --- P1 restart)]
  VM \textbf{Roussette}; overnight batch \texttt{20260516}
  (\file{instruction\_16052026.md}).  First launch cancelled mid-epoch~1;
  second run reached epoch~67/100 then \textbf{VM restart} (admin): only
  encoder trunk on disk (\texttt{e6\_encoder.pt}), MAE decoder not recoverable---
  \textbf{pretrain restarted from scratch} (2026-05-17) with same Hydra preset.
  Prior logs: \texttt{logs/roussette\_e6\_mae\_pretrain\_224\_20260516.log}
  (interrupted), chain \texttt{logs/roussette\_e6\_chain\_20260516.log}.
  \paragraph{Phase 1 (MAE @224, restart).}
  \texttt{expt}=\texttt{track\_a\_ssl\_pretrain\_e6}; honest val excluded from
  pretrain; \texttt{T}=4; \texttt{Sz}=224; \texttt{E}=100; \texttt{wu}=5;
  \texttt{bs}=16; \texttt{lr}=1.5e-4; \texttt{mask}=0.75; \texttt{seed}=42;
  step logs every 100 batches with ISO timestamps (\texttt{log\_timestamps}).
  Log \texttt{logs/roussette\_e6\_mae\_pretrain\_224\_20260517.log};
  PID \texttt{logs/roussette\_e6\_mae\_pretrain\_224\_20260517.pid}.
  \paragraph{Phase 2 (FT @256, auto after P1).}
  \texttt{expt}=\texttt{track\_a\_ssl\_finetune\_e6}; \texttt{Sz}=256;
  \texttt{interpolate\_pos\_embed}=1 (src 224); \texttt{dp}=0.1;
  \texttt{do}=0.0; \texttt{E}=50; \texttt{bs}=12;
  \texttt{tta\_scales}=[0.875, 1.0, 1.125] at submit.
  Watcher \texttt{logs/roussette\_e6\_pt\_then\_ft\_20260517.sh} (PID
  \texttt{logs/roussette\_e6\_pt\_then\_ft\_20260517.pid}) launches FT when
  epoch~100 completes.
  Log \texttt{logs/roussette\_e6\_ft\_champion\_256\_20260517.log}.
\item[\textbf{Run E\_R1 (Sole, 20260517)} (RUNNING)]
  E\_R1 VideoMAE-V2 ViT-B (\texttt{tube\_t}=1, dual masking) then champion FT on
  \textbf{sole.polytechnique.fr} (RTX~3090).
  \paragraph{Phase 1 (DONE).} \texttt{track\_a\_ssl\_pretrain\_e\_r1};
  MAE \texttt{E}=200; encoder \file{e\_r1\_encoder.pt};
  log \texttt{logs/sole\_e\_r1\_chain\_20260517.log} (Phase~1 section).
  \paragraph{Phase 2a (FAILED).} \texttt{track\_a\_ssl\_finetune\_e\_r1},
  \texttt{E}=40; stopped at ep~7/40 (OOM \texttt{Killed}, step~2600/2813).
  Best EMA val top-1 \textbf{33.57\,\%} (ep~6, \file{e\_r1\_ft.pt}).
  \paragraph{Phase 2b (RUNNING).} Resume \texttt{track\_a\_ssl\_finetune\_e\_r1\_resume}
  from \file{e\_r1\_ft.last.pt} (epoch~6); \texttt{E}=60.
  Log \texttt{logs/sole\_e\_r1\_ft\_resume\_20260517.log};
  PID \texttt{logs/sole\_e\_r1\_ft\_resume\_20260517.pid}.
\item[\textbf{Run E3 (Sole, batch 20260516)} (RUNNING)]
  Overnight batch \file{instruction\_16052026.md}: VideoMAE ViT-S at
  \texttt{T}=8 end-to-end (pretrain then champion FT; honest official val).
  Host \textbf{sole.polytechnique.fr} (RTX~3090).
  \paragraph{Phase 1 (DONE).} \texttt{track\_a\_ssl\_pretrain\_e3};
  MAE \texttt{E}=80, \texttt{wu}=5, \texttt{bs}=8, \texttt{grad\_accum}=2,
  \texttt{lr}=1.5e-4, \texttt{mask}=0.75, \texttt{T}=8,
  \texttt{grad\_ckpt}=1; encoder \file{checkpoints/track\_a/ssl/e3\_encoder.pt}
  (symlink to \file{congre\_encoder.pt}, 88\,MB).
  \paragraph{Phase 2 (RUNNING).} \texttt{track\_a\_ssl\_finetune\_e3};
  \texttt{E}=50, \texttt{bs}=4, \texttt{grad\_accum}=2, \texttt{T}=8,
  attentive probe (4 heads); \texttt{tta\_scales}=[0.857, 1.0, 1.143];
  \file{e3\_ft.pt}.  Log \texttt{logs/sole\_e3\_ft\_champion\_t8\_20260516.log};
  PID \texttt{logs/sole\_e3\_ft\_champion\_t8\_20260516.pid}.
  Relaunch after champion YAML fixes (\texttt{training.device}, \texttt{videomix\_alpha}).
  Best val metrics: pending.
\end{description}
\end{minipage}
\end{table}

\subsection{Supervised VideoMAE-style ViT from scratch: verdict}
\label{sec:tracka:vit-from-scratch-verdict}

Every supervised \texttt{video\_mae\_vit} run above---ViT-B vs ViT-S, mean pool
vs attentive probe, with or without class boosting / balancing---stayed near
\(\sim 1/33\) random chance scaled by a small factor (\(\sim 5\)--\(8\,\%\)
train top-1, \(\sim 3\)--\(6\,\%\) val top-1).  The attentive probe (Run~A4)
did \emph{not} unlock learning: mid-training logs still show loss \(\approx
\ln(33)\) and \(\sim 6.5\,\%\) train accuracy, indistinguishable from mean pool
(Run~A2).  Class tricks (rouget / roussette / raie) marginally move numbers
but remain \(\ll\) the \(\sim 39\,\%\) val regime of the dual-stream CNN
(Runs~A1/A1b).  \textbf{Conclusion:} training this ViT \emph{from random init}
on \(\sim 45\,k\) labelled clips alone is not viable for Track~A; the
practical path is \textbf{VideoMAE self-supervised pretraining} on
provided frames (\texttt{track\_a\_videomae\_pretrain}) followed by supervised
fine-tune (\texttt{track\_a\_videomae\_finetune} with
\texttt{model.init\_from}), as configured in this repository.

The overnight \textbf{SSL fish-matrix} (slots 8--16, \file{SSL.md}) implements
this path per codename.  Slot~10 (\textbf{roussette}) completed end-to-end with
best EMA val \textbf{33.61\,\%} (seed/batch drift vs YAML---see
\file{experiments.md}).  Run~A5 (\textbf{truite}, slot~9) finished SSL pretrain only;
slot~12 (\textbf{raie}) best EMA val \textbf{34.80\,\%} (6-head probe);
slot~15 (\textbf{piranha}) best EMA val \textbf{35.60\,\%} (40 FT epochs);
slot~11 (\textbf{rouget}) completed end-to-end with honest val \(\sim 35\,\%\),
public test \(\sim 36\,\%\) (Tab.~\ref{tab:tracka-runlog}).
A follow-up that trained on \texttt{val} clips as well reported \(\sim 60\,\%\)
val but \emph{worse} generalisation---see \file{experiments.md} and the rouget
row above.  Slot~8 (\textbf{barbeau}) was not run.
